% Fragmento público generado desde propietarios íntegros.
% Residencia: 52.13.1.
% Fuente: ../incoming/radion_monografia_20260827/manuscrito/sections/13_sintesis.tex:81-134
\subsubsection{Definición operatoria del radión}

\begin{definition}[Radión HMT--MD]
El radión es la sección orientada de fase, acción y memoria que cumple
simultáneamente:
\begin{enumerate}
\item recorre una unidad de fase \(\theta=1\);
\item barre acción firmada \(\sigma\hret\);
\item pertenece a una vuelta \(T_+=2\pih\) generada coinductivamente;
\item conserva la hoja \(\sigma\), el retorno nonádico y el cilindro de
refinamiento;
\item satisface el empalme exacto
\(1=S_E-\Phi_T+\epsone\).
\end{enumerate}
\end{definition}

El radián convencional es la imagen del radión bajo el functor de olvido
\[
\mathfrak R_\sigma
\longmapsto
\theta=1\in\R/2\pi\Z.
\]
Ese olvido es útil y legítimo para la trigonometría euclídea. Se vuelve
insuficiente cuando la física necesita distinguir acción, signo, retorno,
campo e historia.

\paragraph{Genealogía de la unidad angular natural}

La naturalidad del radián suele justificarse porque la derivada de
\(e^{i\theta}\), \(\sin\theta\) y \(\cos\theta\) adopta su forma más simple
cuando \(\theta\) se mide como cociente arco/radio. HMT conserva esa ventaja,
pero la explica desde una capa anterior:
\[
\theta\quad\leftrightarrow\quad
\frac{\mathcal A(\theta)}{\hret}.
\]
La fase es natural porque mide acción en unidades de \(\hret\); la vuelta es
natural porque encierra \(\hfull\). El cociente arco/radio es la publicación
euclídea de una unidad que ya posee estructura simpléctica y memoria.

\paragraph{Compatibilidad entre elipse de acción y crecimiento dimensional}

El círculo es la forma que resulta al olvidar la anisotropía:
\(C^*=0\Rightarrow\eta=0\Rightarrow a=b\). La elipse es más informativa porque
conserva dos autovalores, dos ejes y una orientación de intercambio a producto
fijo. En ese sentido, la elipse es una prolongación del círculo, no un círculo
deformado por accidente.

La profundidad de Klein y el lift helicoidal añaden dos tipos de información
que el contorno plano no contiene: distancia interior y memoria de retorno.
Ésta es la imagen ontológica de Mecánica Dimensional que emerge del radión:
una dimensión nueva aparece cuando se exige conservar un dato independiente
que la proyección anterior debía olvidar.

